% Part: incompleteness
% Chapter: incompleteness-provability
% Section: introduction

\documentclass[../../../include/open-logic-section]{subfiles}

\begin{document}

\olfileid{inc}{inp}{int}

\olsection{పరిచయం}

సహజ సంఖ్యల వంటి గణిత నిర్మాణానికి సంబంధించిన స్వీకృతాల వ్యవస్థ, ఆ నిర్మాణం
గురించిన అన్ని సత్య ప్రకటనలనూ నిగమించేందుకు వీలు కల్పించనంతవరకు సరిపోదని
హిల్బర్ట్ భావించాడు. ఆ తర్వాత ఆయనకు అధికారిక నిగమన వ్యవస్థలపై కలిగిన ఆసక్తితో
దీనిని కలిపి చూస్తే, సహజ సంఖ్యల గురించి తర్కించేందుకు మనం వాడే అధికారిక
వ్యవస్థలు అవైరుధ్యంగా మాత్రమే కాకుండా \emph{సంపూర్ణంగా} కూడా ఉండాలని---అంటే,
వాటి భాషలోని ప్రతి ప్రకటన గాని దాని నిషేధం గాని
\tetoken{వ్యుత్పాదించదగినదిగా}{!!{derivable}} ఉండాలని---మనం హామీ ఇవ్వాలి అని ఆయన భావించినట్లు
తెలుస్తుంది. అలాంటి స్వీకృతాల వ్యవస్థ ఏదీ లేదని గ్యోడెల్ మొదటి అసంపూర్ణతా
సిద్ధాంతం చూపుతుంది: అంకగణితానికి సంపూర్ణమైన, అవైరుధ్యమైన,
\tetoken{స్వీకృతీకరించదగిన}{!!{axiomatizable}} అధికారిక వ్యవస్థ ఏదీ లేదు. నిజానికి,
``తగినంత బలమైన,'' అవైరుధ్యమైన, \tetoken{స్వీకృతీకరించదగిన}{!!{axiomatizable}} గణిత సిద్ధాంతం
ఏదీ సంపూర్ణం కాదు.

ఆధునిక (``సాంప్రదాయిక'') గణితానికి సమర్థన కల్పించే తన కార్యక్రమానికి
కేంద్రబిందువైన హిల్బర్ట్ మరొక లక్ష్యం ఇంకా ముఖ్యమైనది: సాంప్రదాయిక తర్కాన్ని
ప్రాతినిధ్యం చేసే అధికారిక వ్యవస్థలకు పరిసమాప్త అవైరుధ్య నిరూపణలను కనుగొనడం.
కాబట్టి హిల్బర్ట్ కార్యక్రమం దృష్టిలో గ్యోడెల్ రెండవ అసంపూర్ణతా సిద్ధాంతం ఇంకా
పెద్ద దెబ్బ. మొదటి సిద్ధాంతంలాగే రెండవ అసంపూర్ణతా సిద్ధాంతాన్నీ కొంత అస్పష్టంగా
చెప్పవచ్చు. స్థూలంగా, తగినంత బలమైన అంకగణిత సిద్ధాంతం తన స్వంత అవైరుధ్యాన్ని
నిరూపించలేదని అది చెబుతుంది. ఇక్కడ ``తగినంత బలమైనది'' అనేది~$\Th{Q}$ కంటే
కొద్దిగా ఎక్కువను కలిగి ఉండాలి.

అసంపూర్ణతా సిద్ధాంతానికి గ్యోడెల్ ఇచ్చిన అసలు నిరూపణ వెనుక ఆలోచన ఎపిమెనిడీస్
విరోధాభాసంలో కనిపిస్తుంది. క్రీట్‌కు చెందిన ఎపిమెనిడీస్, క్రీట్‌వాసులందరూ
అబద్ధికులని అన్నాడు; ఆ విరోధాభాసానికి మరింత ప్రత్యక్ష రూపం ``ఈ వాక్యం
అసత్యం'' అనే ప్రకటన. సత్యం స్థానంలో \tetoken{వ్యుత్పాద్యతను}{!!{derivability}} పెట్టడం ద్వారా,
తానే \tetoken{వ్యుత్పాదించదగినది}{!!{derivable}} కాదని పరోక్షంగా చెప్పే
\tetoken{ఒక వాక్యాన్ని}{!!a{sentence}} గ్యోడెల్ అధికారికీకరించగలిగాడు. ఆ
\tetoken{వాక్యం}{!!{sentence}} \tetoken{వ్యుత్పాదించదగినదైతే}{!!{derivable}}, సిద్ధాంతం వైరుధ్యమవుతుంది.
తాను పరిశీలించిన స్వీకృతాల వ్యవస్థ నుంచి ఆ \tetoken{వాక్యం}{!!{sentence}} యొక్క నిషేధం కూడా
\tetoken{వ్యుత్పాదించదగినది}{!!{derivable}} కాదని గ్యోడెల్ చూపించాడు. (ఈ రెండవ భాగానికి
సిద్ధాంతం~$\Th{T}$ ``$\omega$-అవైరుధ్యమైనది'' అని గ్యోడెల్ భావించాల్సి వచ్చింది.
$\omega$-అవైరుధ్యం అవైరుధ్యానికి సంబంధించినదే గాని దానికంటే బలమైన
లక్షణం.\footnote{అంటే ప్రతి $\omega$-అవైరుధ్య సిద్ధాంతమూ అవైరుధ్యమైనదే;
దీని విపర్యయం మాత్రం నిజం కాదు.} గ్యోడెల్ తర్వాత కొన్ని సంవత్సరాలకు,
సాధారణంగా~$\Th{T}$ అవైరుధ్యమైనదని మాత్రమే ఊహిస్తే చాలునని రాసర్ చూపించాడు.)

తననే సూచించే \tetoken{ఒక వాక్యాన్ని}{!!a{sentence}} ఎలా నిర్మించవచ్చో అర్థం చేసుకోవడం మొదటి
సవాలు. $\Th{Q}$ భాషలోని ప్రతి \tetoken{సూత్రం}{!!{formula}}~$!A$కు, $\Gn{!A}$కు
అనుగుణమైన సంఖ్యాంకాన్ని~$\gn{!A}$ సూచించనివ్వండి. దీని అర్థాన్ని జాగ్రత్తగా
చూడండి: $!A$ అనేది $\Th{Q}$ భాషలోని \tetoken{ఒక సూత్రం}{!!a{formula}}, $\Gn{!A}$ ఒక సహజ
సంఖ్య, $\gn{!A}$ మాత్రం $\Th{Q}$ భాషలోని ఒక \emph{పదం}. అందువల్ల
$\Th{Q}$ భాషలోని ప్రతి \tetoken{సూత్రం}{!!{formula}}~$!A$కు అదే భాషలోని పదమైన
\emph{పేరు}~$\gn{!A}$ ఉంటుంది; అది $\Th{Q}$ భాషలోని పదమే; దీని వల్ల $\Th{Q}$ భాషలోని
\tetoken{సూత్రాలు}{!!{formula}s} ఇతర \tetoken{సూత్రాల}{!!{formula}s} గురించి విషయాలు ``చెప్పగలిగే'' భావనాత్మక
చట్రం మనకు లభిస్తుంది. కింది ఉపసిద్ధాంతాన్ని స్థిరబిందు ఉపసిద్ధాంతం అంటారు.

\begin{lem}
$\Th{Q}$ను విస్తరించే ఏ సిద్ధాంతమైనా $\Th{T}$గా, $x$ మాత్రమే స్వేచ్ఛగా ఉన్న
ఏ \tetoken{సూత్రమైనా}{!!{formula}} $!B(x)$గా తీసుకోండి. అప్పుడు $\Th{T} \Proves!A \liff
!B(\gn{!A})$ అయ్యే \tetoken{ఒక వాక్యం}{!!a{sentence}}~$!A$ ఉంటుంది.
\end{lem}

ఏ లక్షణం $!B(x)$ ఇచ్చినా, ``$!B(x)$ నాకు నిజం'' అని ప్రకటించే
\tetoken{ఒక వాక్యం}{!!a{sentence}}~$!A$ ఉంటుందని, ఈ విషయాన్ని $\Th{T}$ ``తెలుసుకుంటుందని'' ఈ
ఉపసిద్ధాంతం చెబుతుంది.

అలాంటి \tetoken{ఒక వాక్యాన్ని}{!!a{sentence}} ఎలా నిర్మించాలి? క్వైన్‌కు చెందిన, ఎపిమెనిడీస్
విరోధాభాసపు ఈ రూపాన్ని పరిశీలించండి:
\begin{quote}
``తన ఉల్లేఖనాన్ని ముందు ఉంచితే అసత్యాన్ని ఇస్తుంది'' అనే పదబంధం, తన
ఉల్లేఖనాన్ని ముందు ఉంచితే అసత్యాన్ని ఇస్తుంది.
\end{quote}
ఈ వాక్యం నేరుగా తనను తాను సూచించదు. ఉల్లేఖన చిహ్నాల మధ్యనున్న వాక్యనిర్మాణ
వస్తువుల గురించి మాత్రమే అది ఒక ప్రకటన చేస్తుంది; ఆ విషయంలో ఇది కింది
వాక్యాలతో సమానమే:
\begin{enumerate}
\item ``రాబర్ట్'' మంచి పేరు.
\item ``నేను పరుగెత్తాను.'' చిన్న వాక్యం.
\item ``మూడు పదాలు ఉన్నాయి'' అనే పదబంధంలో మూడు పదాలు ఉన్నాయి.
\end{enumerate}
కానీ ``తన ఉల్లేఖనాన్ని ముందు ఉంచితే అసత్యాన్ని ఇస్తుంది'' అనే పదబంధానికి,
దాని స్వంత ఉల్లేఖిత రూపాన్ని ముందు ఉంచితే ఏమవుతుంది? మనకు అసలు వాక్యమే
వస్తుంది!{} సంక్షిప్తంగా, ఆ వాక్యం తాను అసత్యమని ప్రకటిస్తుంది.

\end{document}
